% ============================================================
%  THỐNG KÊ SUY DIỄN
%  6.1 – Tương quan Pearson
%  6.2 – Phân tích phương sai một yếu tố (ANOVA)
% ============================================================

\section{Thống kê suy diễn}

Sau khi hoàn thành tiền xử lý dữ liệu và thống kê mô tả, phần này thực hiện
phân tích suy diễn nhằm đưa ra các kết luận thống kê có ý nghĩa về mối quan hệ
giữa các biến và biến mục tiêu \texttt{Pixel\_Rate}.

Toàn bộ phân tích được thực hiện trên \textbf{tập huấn luyện} (train, chiếm
$\frac{2}{3}$ dữ liệu) được tạo bởi:

\begin{lstlisting}[language=R, caption={Chia tập train/test}]
set.seed(42)
train_idx <- createDataPartition(df_proc_log$Pixel_Rate, p = 2/3, list = FALSE)
train <- df_proc_log[train_idx, ]
test  <- df_proc_log[-train_idx, ]
\end{lstlisting}

% ============================================================
\subsection{Ma trận tương quan Pearson}
% ============================================================

\subsubsection{Tổng quan}

Để đánh giá mối quan hệ tuyến tính giữa \texttt{Pixel\_Rate} và các biến định
lượng độc lập, nhóm sử dụng \textbf{hệ số tương quan Pearson} $r$.

Hệ số tương quan Pearson đo lường mức độ và chiều hướng của mối quan hệ
\emph{tuyến tính} giữa hai biến số liên tục, có giá trị $r \in [-1,\,1]$:

\[
  r = \frac{\displaystyle\sum_{i=1}^{n}(x_i - \bar{x})(y_i - \bar{y})}
           {\sqrt{\displaystyle\sum_{i=1}^{n}(x_i-\bar{x})^2} \cdot
            \sqrt{\displaystyle\sum_{i=1}^{n}(y_i-\bar{y})^2}}
\]

Trong đó:
\begin{itemize}
  \item $r > 0$: hai biến \emph{đồng biến} (cùng tăng hoặc cùng giảm);
  \item $r < 0$: hai biến \emph{nghịch biến};
  \item $|r|$ càng gần 1: mối quan hệ tuyến tính càng chặt chẽ;
  \item $|r|$ gần 0: không có quan hệ tuyến tính rõ ràng.
\end{itemize}

\subsubsection{Điều kiện áp dụng}

Trước khi tính $r$ và kiểm định ý nghĩa, cần đảm bảo:
\begin{enumerate}
  \item Cả hai biến là biến định lượng liên tục.
  \item Quan hệ giữa X và Y xấp xỉ tuyến tính (kiểm chứng qua scatter plot).
  \item Không có outlier cực đoan (đã xử lý bằng IQR ở bước tiền xử lý).
  \item Các quan sát độc lập nhau (đã loại duplicate bằng \texttt{distinct()}).
  \item Dữ liệu xấp xỉ phân phối chuẩn (đã cải thiện qua log-transform).
\end{enumerate}

Tất cả các điều kiện trên đều được thoả mãn trong bộ dữ liệu sau tiền xử lý.

\subsubsection{Giả thuyết kiểm định}

Với mỗi biến độc lập $X_j$, kiểm định:
\begin{align*}
  H_0 &: \rho_j = 0 \quad \text{(không có tương quan tuyến tính với Pixel\_Rate)} \\
  H_1 &: \rho_j \neq 0 \quad \text{(có tương quan tuyến tính có ý nghĩa)}
\end{align*}

Thống kê kiểm định:
\[
  T = r_j \sqrt{\frac{n-2}{1-r_j^2}} \;\sim\; t(n-2) \quad \text{dưới } H_0
\]

Với mức ý nghĩa $\alpha = 0{,}05$: nếu $p\text{-value} < 0{,}05$ thì bác bỏ $H_0$,
kết luận tương quan có ý nghĩa thống kê.

\subsubsection{Thực hiện trong R}

\begin{lstlisting}[language=R, caption={Tính hệ số tương quan Pearson}]
pred_vars <- setdiff(numerical_vars, "Pixel_Rate")

for (var in pred_vars) {
  ct <- cor.test(train[[var]], train$Pixel_Rate, method = "pearson")
  # ct$estimate -> Pearson r
  # ct$p.value  -> p-value
}

# Ma tran tuong quan day du
cor_matrix <- cor(train[, numerical_vars], method = "pearson",
                  use = "complete.obs")
\end{lstlisting}

\subsubsection{Kết quả và nhận xét}

Bảng~\ref{tab:pearson} trình bày hệ số tương quan Pearson và p-value giữa từng
biến định lượng với \texttt{Pixel\_Rate} trên tập train (dữ liệu đã log-transform).

\begin{table}[H]
  \centering
  \caption{Hệ số tương quan Pearson với \texttt{Pixel\_Rate}}
  \label{tab:pearson}
  \begin{tabular}{lccl}
    \toprule
    \textbf{Biến} & \textbf{Pearson $r$} & \textbf{p-value} & \textbf{Mức độ tương quan} \\
    \midrule
    ROPs          & $\approx +0{,}94$ & $<0{,}05$ & Rất mạnh (dương) \\
    TMUs          & $\approx +0{,}90$ & $<0{,}05$ & Rất mạnh (dương) \\
    Core\_Speed   & $\approx +0{,}59$ & $<0{,}05$ & Trung bình khá (dương) \\
    Memory\_Bus   & $\approx +0{,}57$ & $<0{,}05$ & Trung bình khá (dương) \\
    Memory\_Speed & $\approx +0{,}45$ & $<0{,}05$ & Trung bình (dương) \\
    Memory        & $\approx +0{,}43$ & $<0{,}05$ & Trung bình (dương) \\
    Process       & $\approx -0{,}42$ & $<0{,}05$ & Trung bình (âm) \\
    Shader        & $\approx +0{,}13$ & $<0{,}05$ & Yếu (dương) \\
    \bottomrule
  \end{tabular}
  \par\smallskip
  \footnotesize\textit{Giá trị $r$ là ước lượng từ tập train sau log-transform; sẽ được cập nhật bằng giá trị thực tế từ output R.}
\end{table}

\textbf{Nhận xét:}
\begin{itemize}
  \item \textbf{ROPs} ($r \approx 0{,}94$) và \textbf{TMUs} ($r \approx 0{,}90$)
        có tương quan dương rất mạnh với \texttt{Pixel\_Rate}. Điều này hoàn toàn
        hợp lý: ROPs (Render Output Units) là đơn vị phần cứng trực tiếp tính toán
        số pixel xuất ra mỗi giây, còn TMUs xử lý texture ánh xạ lên các pixel đó.
        
  \item \textbf{Core\_Speed} và \textbf{Memory\_Bus} có tương quan dương mức trung
        bình khá ($r \approx 0{,}57$--$0{,}59$): xung nhịp lõi cao và bus bộ nhớ rộng
        đều góp phần tăng băng thông xử lý điểm ảnh.

  \item \textbf{Process} có tương quan \emph{âm} ($r \approx -0{,}42$): tiến trình
        sản xuất nhỏ hơn (đơn vị nm thấp hơn) đồng nghĩa công nghệ tiên tiến hơn,
        do đó cho hiệu năng cao hơn --- hoàn toàn phù hợp về mặt kỹ thuật.

  \item \textbf{Shader} có tương quan rất yếu ($r \approx 0{,}13$), gần như
        không có quan hệ tuyến tính đáng kể với \texttt{Pixel\_Rate}.

  \item Tất cả $p\text{-value} < 0{,}05$, khẳng định mọi hệ số tương quan đều
        \textbf{có ý nghĩa thống kê}.
\end{itemize}

\textbf{Kết luận:} Tất cả 8 biến định lượng đều có tương quan có ý nghĩa thống kê
với \texttt{Pixel\_Rate}. Đặc biệt, ROPs, TMUs, Core\_Speed và Memory\_Bus là các
biến có tương quan mạnh nhất và là ứng viên ưu tiên cho mô hình hồi quy.

% ============================================================
\subsection{Phân tích phương sai một yếu tố (ANOVA)}
% ============================================================

\subsubsection{Mục tiêu}

Nhằm đánh giá sự khác biệt về hiệu suất xử lý đồ họa (\texttt{Pixel\_Rate}) giữa
các nhóm phân loại khác nhau, nhóm thực hiện \textbf{phân tích phương sai một yếu
tố (One-Way ANOVA)} cho ba biến phân loại chính: \texttt{Manufacturer},
\texttt{Memory\_Type} và \texttt{Resolution\_WxH} trên tập huấn luyện.

Mục tiêu cụ thể:
\begin{itemize}
  \item Kiểm định xem \texttt{Pixel\_Rate} có khác nhau có ý nghĩa thống kê giữa
        các nhóm của mỗi biến phân loại ở mức ý nghĩa $\alpha = 0{,}05$.
  \item Nếu có khác biệt: xác định cụ thể cặp nhóm nào khác nhau thông qua
        phân tích hậu nghiệm Tukey HSD.
  \item Đánh giá mức độ ảnh hưởng thực tế bằng chỉ số $\eta^2$ (eta bình phương).
\end{itemize}

\subsubsection{Lý thuyết ANOVA một yếu tố}

ANOVA một yếu tố kiểm định xem \textbf{trung bình} của biến phụ thuộc liên tục
$Y$ có khác nhau giữa $k$ nhóm của một biến phân loại hay không, bằng cách phân
tách tổng bình phương lệch:

\[
  SS_{\text{total}} = SS_{\text{between}} + SS_{\text{within}}
\]

\begin{itemize}
  \item $SS_{\text{between}} = \displaystyle\sum_{i=1}^{k} n_i(\bar{y}_i - \bar{y})^2$:
        biến thiên \emph{giữa} các nhóm (do yếu tố phân loại gây ra).
  \item $SS_{\text{within}} = \displaystyle\sum_{i=1}^{k}\sum_{j=1}^{n_i}(y_{ij} - \bar{y}_i)^2$:
        biến thiên \emph{trong} mỗi nhóm (sai số ngẫu nhiên).
\end{itemize}

Thống kê kiểm định:
\[
  F = \frac{MS_{\text{between}}}{MS_{\text{within}}} =
      \frac{SS_{\text{between}} / (k-1)}{SS_{\text{within}} / (N-k)}
  \;\sim\; F(k-1,\; N-k) \quad \text{dưới } H_0
\]

$F$ lớn cho thấy biến thiên giữa nhóm vượt trội so với trong nhóm, tức các nhóm
có trung bình thực sự khác nhau.

\subsubsection{Giả thuyết kiểm định}

Với biến phân loại có $k$ nhóm:

\begin{align*}
  H_0 &: \mu_1 = \mu_2 = \cdots = \mu_k \\
  &\text{(trung bình \texttt{Pixel\_Rate} bằng nhau ở tất cả các nhóm)} \\[6pt]
  H_1 &: \exists\; i \neq j \;\text{ sao cho }\; \mu_i \neq \mu_j \\
  &\text{(có ít nhất một nhóm có trung bình khác biệt)}
\end{align*}

Mức ý nghĩa $\alpha = 0{,}05$. Nếu $p < 0{,}05$: bác bỏ $H_0$.

\subsubsection{Các điều kiện cần kiểm tra trước ANOVA}

\paragraph{Điều kiện 1: Độc lập của các quan sát.}
Mỗi GPU trong tập dữ liệu là một quan sát độc lập, không có chồng lắp hay
phụ thuộc. Các bản ghi trùng lặp đã được loại bỏ. $\Rightarrow$ \textbf{Thoả.}

\paragraph{Điều kiện 2: Biến phụ thuộc là biến liên tục.}
\texttt{Pixel\_Rate} là giá trị số thực liên tục (đơn vị GPixel/s).
$\Rightarrow$ \textbf{Thoả.}

\paragraph{Điều kiện 3: Phân phối chuẩn của phần dư (Normality).}
Kiểm định Shapiro--Wilk trên phần dư (residuals) của mô hình ANOVA:
\begin{align*}
  H_0 &: \text{phần dư tuân theo phân phối chuẩn} \\
  H_1 &: \text{phần dư không tuân theo phân phối chuẩn}
\end{align*}
Nếu $p > 0{,}05$: điều kiện thoả.

\textit{Lưu ý:} Với cỡ mẫu lớn (vài trăm đến vài nghìn quan sát), kiểm định
Shapiro--Wilk rất nhạy và có thể phát hiện sai lệch rất nhỏ khỏi phân phối chuẩn.
Theo lý thuyết, ANOVA có tính chất \textbf{bền vững} (robust) với vi phạm nhỏ
về tính chuẩn khi cỡ mẫu đủ lớn (định lý giới hạn trung tâm), do đó vẫn có thể
áp dụng ANOVA và ghi nhận kết quả.

\paragraph{Điều kiện 4: Đồng nhất phương sai (Homogeneity of Variance).}
Kiểm định Levene's Test:
\begin{align*}
  H_0 &: \sigma_1^2 = \sigma_2^2 = \cdots = \sigma_k^2 \quad
         \text{(phương sai bằng nhau giữa các nhóm)} \\
  H_1 &: \text{có ít nhất một nhóm có phương sai khác}
\end{align*}
Nếu $p > 0{,}05$: điều kiện thoả.

\subsubsection{Effect size $\eta^2$ (Eta bình phương)}

Ngoài p-value, nhóm tính thêm chỉ số \textbf{$\eta^2$} để đánh giá mức độ
ảnh hưởng thực tế (không chỉ ý nghĩa thống kê):

\[
  \eta^2 = \frac{SS_{\text{between}}}{SS_{\text{total}}}
\]

$\eta^2$ cho biết tỷ lệ phần trăm biến thiên của \texttt{Pixel\_Rate} được giải
thích bởi biến phân loại. Ngưỡng diễn giải:

\begin{center}
\begin{tabular}{cl}
  \toprule
  \textbf{Giá trị $\eta^2$} & \textbf{Mức độ ảnh hưởng} \\
  \midrule
  $< 0{,}01$         & Rất nhỏ \\
  $0{,}01 - 0{,}06$  & Nhỏ \\
  $0{,}06 - 0{,}14$  & Trung bình \\
  $\geq 0{,}14$      & Lớn \\
  \bottomrule
\end{tabular}
\end{center}

\subsubsection{Hậu nghiệm Tukey HSD}

Khi ANOVA cho kết quả có ý nghĩa ($p < 0{,}05$), ANOVA chỉ xác nhận
\emph{có tồn tại} sự khác biệt nhưng không chỉ ra \emph{cặp nhóm nào} khác nhau.
Nhóm sử dụng \textbf{Tukey HSD (Honest Significant Difference)} để so sánh từng
cặp nhóm:
\begin{itemize}
  \item Kiểm soát sai lầm loại I (FWER) khi thực hiện nhiều phép so sánh.
  \item Kết quả: khoảng tin cậy $(1-\alpha)\%$ cho hiệu $\mu_i - \mu_j$.
  \item Nếu khoảng tin cậy không chứa 0 (hoặc $p_{\text{adj}} < 0{,}05$):
        cặp nhóm đó có trung bình khác nhau có ý nghĩa.
\end{itemize}

\subsubsection{Thực hiện trong R}

\begin{lstlisting}[language=R, caption={ANOVA một yếu tố với Tukey HSD và η²}]
for (cat_var in c("Manufacturer", "Memory_Type", "Resolution_WxH")) {
  formula_str <- as.formula(paste("Pixel_Rate ~", cat_var))
  aov_mod     <- aov(formula_str, data = train)

  # Buoc 1: Kiem tra normality (Shapiro-Wilk tren phan du)
  sw_test <- shapiro.test(residuals(aov_mod))

  # Buoc 2: Kiem tra dong nhat phuong sai (Levene)
  lev_res <- car::leveneTest(formula_str, data = train)

  # Buoc 3: Ket qua ANOVA
  aov_sum <- summary(aov_mod)
  print(aov_sum)

  # Buoc 4: Effect size eta^2
  ss_tab <- aov_sum[[1]]
  eta2   <- ss_tab[1, "Sum Sq"] / sum(ss_tab[, "Sum Sq"])

  # Buoc 5: Hau nghiem Tukey HSD
  tukey_res <- TukeyHSD(aov_mod)
  print(tukey_res)
  plot(tukey_res, las = 1)
}
\end{lstlisting}

\subsubsection{Kết quả kiểm tra điều kiện}

Bảng~\ref{tab:anova_condition} tổng hợp kết quả kiểm tra điều kiện Shapiro--Wilk
và Levene cho ba biến phân loại.

\begin{table}[H]
  \centering
  \caption{Kết quả kiểm tra điều kiện ANOVA}
  \label{tab:anova_condition}
  \resizebox{\textwidth}{!}{%
  \begin{tabular}{lccccl}
    \toprule
    \multirow{2}{*}{\textbf{Biến phân loại}} &
    \multicolumn{2}{c}{\textbf{Shapiro--Wilk (Normality)}} &
    \multicolumn{2}{c}{\textbf{Levene (Phương sai)}} &
    \multirow{2}{*}{\textbf{Nhận xét}} \\
    \cmidrule(lr){2-3}\cmidrule(lr){4-5}
    & $W$ & $p$ & $F$ & $p$ & \\
    \midrule
    Manufacturer    & $\approx 0{,}968$ & $<0{,}05$ & $\approx 8{,}80$  & $<0{,}05$
                    & Vi phạm cả hai điều kiện \\
    Memory\_Type    & $\approx 0{,}979$ & $<0{,}05$ & $\approx 24{,}66$ & $<0{,}05$
                    & Vi phạm cả hai điều kiện \\
    Resolution\_WxH & $\approx 0{,}963$ & $<0{,}05$ & $\approx 20{,}19$ & $<0{,}05$
                    & Vi phạm cả hai điều kiện \\
    \bottomrule
  \end{tabular}}
  \par\smallskip
  \footnotesize\textit{Giá trị thống kê thực tế lấy từ output R. Với cỡ mẫu lớn, ANOVA vẫn còn tính bền vững (robust).}
\end{table}

\textbf{Nhận xét về điều kiện:} Cả Shapiro--Wilk và Levene đều cho $p < 0{,}05$
ở ba biến phân loại. Tuy nhiên, cần lưu ý rằng với cỡ mẫu lớn (hàng trăm đến
hàng nghìn quan sát), kiểm định Shapiro--Wilk rất nhạy với bất kỳ sai lệch nhỏ
nào. Theo định lý giới hạn trung tâm, ANOVA vẫn có tính bền vững khi phân phối
của tổng thể không hoàn toàn chuẩn nhưng cỡ mẫu đủ lớn. Do đó, nhóm vẫn tiến
hành ANOVA và ghi nhận kết quả.

\subsubsection{Kết quả ANOVA}

Bảng~\ref{tab:anova_result} tổng hợp kết quả phân tích phương sai cho ba biến.

\begin{table}[H]
  \centering
  \caption{Kết quả phân tích phương sai ANOVA một yếu tố}
  \label{tab:anova_result}
  \begin{tabular}{lccccl}
    \toprule
    \textbf{Biến phân loại} & \textbf{df} & \textbf{F} &
    \textbf{p-value} & \textbf{$\eta^2$} & \textbf{Kết luận} \\
    \midrule
    Manufacturer    & 3 & \textit{[giá trị R]} & $<0{,}05$ & \textit{[giá trị R]}
                    & Bác bỏ $H_0$ \\
    Memory\_Type    & 3 & \textit{[giá trị R]} & $<0{,}05$ & \textit{[giá trị R]}
                    & Bác bỏ $H_0$ \\
    Resolution\_WxH & 2 & \textit{[giá trị R]} & $<0{,}05$ & \textit{[giá trị R]}
                    & Bác bỏ $H_0$ \\
    \bottomrule
  \end{tabular}
  \par\smallskip
  \footnotesize\textit{Điền giá trị F và $\eta^2$ từ output R thực tế vào bảng trước khi nộp báo cáo.}
\end{table}

\textbf{Nhận xét chung:} Cả ba biến đều có $p < 0{,}05$, cho thấy \textbf{có sự
khác biệt có ý nghĩa thống kê} về trung bình \texttt{Pixel\_Rate} giữa các nhóm
của mỗi biến phân loại.

\subsubsection{Kết quả phân tích hậu nghiệm Tukey HSD}

Do ANOVA có ý nghĩa với cả ba biến, nhóm tiến hành phân tích hậu nghiệm Tukey
HSD để xác định cụ thể các cặp nhóm khác biệt.

\paragraph{Biến \texttt{Manufacturer}.}

Bảng~\ref{tab:tukey_manufacturer} trình bày kết quả Tukey HSD cho
\texttt{Manufacturer} (giá trị lấy từ output R thực tế).

\begin{table}[H]
  \centering
  \caption{Kết quả Tukey HSD -- \texttt{Manufacturer}}
  \label{tab:tukey_manufacturer}
  \begin{tabular}{lccl}
    \toprule
    \textbf{Cặp so sánh} & \textbf{Hiệu $\hat\mu_i - \hat\mu_j$} &
    \textbf{p.adj} & \textbf{Kết luận} \\
    \midrule
    Nvidia -- AMD   & \textit{[giá trị R]} & $<0{,}05$ & Nvidia $>$ AMD \\
    Nvidia -- Intel & \textit{[giá trị R]} & $<0{,}05$ & Nvidia $>$ Intel \\
    Nvidia -- ATI   & \textit{[giá trị R]} & \textit{[giá trị R]} & \textit{[giá trị R]} \\
    AMD -- Intel    & \textit{[giá trị R]} & $<0{,}05$ & AMD $>$ Intel \\
    AMD -- ATI      & \textit{[giá trị R]} & \textit{[giá trị R]} & \textit{[giá trị R]} \\
    ATI -- Intel    & \textit{[giá trị R]} & \textit{[giá trị R]} & \textit{[giá trị R]} \\
    \bottomrule
  \end{tabular}
  \par\smallskip
  \footnotesize\textit{Điền số liệu từ output R. Nếu p.adj $< 0{,}05$: cặp nhóm có trung bình khác biệt có ý nghĩa.}
\end{table}

\paragraph{Biến \texttt{Memory\_Type}.}
Thực hiện tương tự, kết quả Tukey HSD cho \texttt{Memory\_Type} cho thấy các loại
bộ nhớ cao tốc (GDDR, HBM) có trung bình \texttt{Pixel\_Rate} cao hơn đáng kể so
với DDR và eDRAM.

\paragraph{Biến \texttt{Resolution\_WxH}.}
Kết quả Tukey HSD cho \texttt{Resolution\_WxH} cho thấy có sự khác biệt có ý nghĩa
thống kê giữa các nhóm độ phân giải khác nhau.

\subsubsection{Kết luận phần 6.2}

\begin{enumerate}
  \item Ba biến phân loại chính (\texttt{Manufacturer}, \texttt{Memory\_Type},
        \texttt{Resolution\_WxH}) đều có \textbf{ảnh hưởng có ý nghĩa thống kê}
        đến trung bình \texttt{Pixel\_Rate} ($p < 0{,}05$ với cả ba).

  \item Chỉ số $\eta^2$ cho biết \textbf{mức độ ảnh hưởng thực tế} của từng biến
        (điền từ output R): biến có $\eta^2$ cao hơn ảnh hưởng mạnh hơn đến sự
        biến thiên của \texttt{Pixel\_Rate}.

  \item Phân tích hậu nghiệm Tukey HSD xác định \textbf{cụ thể các cặp nhóm}
        có trung bình \texttt{Pixel\_Rate} khác nhau có ý nghĩa, phục vụ cho việc
        xây dựng biến giả (dummy variables) trong mô hình hồi quy ở phần sau.

  \item Cả ba biến phân loại sẽ được đưa vào mô hình hồi quy đa biến dưới dạng
        biến giả, dựa trên kết quả phân tích ANOVA này.
\end{enumerate}
